\documentclass[11pt]{amsart}
\usepackage[T1]{fontenc}
\usepackage{lmodern}
\usepackage[margin=1in]{geometry}
\usepackage{amsmath,amssymb,amsthm,booktabs,array,xcolor}
\usepackage{microtype}
\usepackage[colorlinks=true,linkcolor=blue!45!black,citecolor=blue!45!black,urlcolor=blue!45!black]{hyperref}
\newtheorem{theorem}{Theorem}[section]
\newtheorem{lemma}[theorem]{Lemma}
\newtheorem{proposition}[theorem]{Proposition}
\newtheorem{corollary}[theorem]{Corollary}
\theoremstyle{definition}
\newtheorem{definition}[theorem]{Definition}
\theoremstyle{remark}
\newtheorem{remark}[theorem]{Remark}
\newcommand{\ceil}[1]{\left\lceil#1\right\rceil}
\newcommand{\R}{\mathbb R}
\newcommand{\Amin}{A_{\min}}
\newcommand{\Rsec}{R_{\mathrm{sec}}}
\newcommand{\repo}[2]{\href{https://github.com/shjvandekamp/math/blob/main/topics/cylinder-covering/#1}{#2}}
\hypersetup{pdftitle={A quartic overlap certificate and effective finite cylinder covers of a regular tetrahedron}}
\title[Effective finite cylinder covers of a tetrahedron]{A quartic overlap certificate and effective finite\texorpdfstring{\\}{} cylinder covers of a regular tetrahedron}
\date{15 October 2026}
\begin{document}
\begin{abstract}
We give an explicit finite version of the angular cylinder construction
for a regular tetrahedron, using the secant axes and matching angular
boundaries of the OpenAI family-100 manuscripts. A polynomial
simultaneous-failure identity permits a quartic radial correction
$9e^4/16$. An explicit perturbation bound transfers this certificate to
every finite positive-width angular partition and yields compact,
nondegenerate triangular perpendicular bases covering the entire closed
tetrahedron. For a uniform partition into $N$ intervals, the normalized
area differs from its limiting integral by at most $2e/N$ for
$0<e\le1/10$. In particular, $N\ge\ceil{480/e}$ guarantees a saving
of at least $e^2/240$ below one half. On the schedule
$N=\ceil{2/e^2}$ the actual area is
$1/2-e^2/60+e^3/2+O(e^4)$, with an explicit quartic error bound.
An exact rational certificate places a 20,000-cylinder example in
$[0.499921585997001674,\,0.499921585997001774]$ times the minimum
projection area. The comparison is with identified published recipes on
their explicit sufficient ranges and prescribed count conventions;
no global optimality or domination of every possible generic-method
instance is asserted.
\end{abstract}
\maketitle

\section{Scope, geometry, and the main theorem}

The existence of a finite triangular-cylinder cover below half the minimum
projection area is already established in the OpenAI family-100
manuscripts~\cite{Angular,Radial,Ruled,Slope}. This note uses the angular
paper's secant directions, shared sector boundaries, and first-crossing
selection~\cite{Angular}. Its contribution is the quartic overlap identity
in Lemma~\ref{lem:square} and the explicit finite perturbation estimate
that follows it. These change the proved cap margins and parameter
ranges within an existing geometric construction. In particular, neither
the secant sectors themselves nor the original counterexample to the
half-area conjecture are claimed as new here.

All four source manuscripts and the cited formal source files are pinned
to the OpenAI repository revision
\texttt{1581d4cbdd69a0faf34b2089614fe2a814316a6a}. The public
\repo{sources/manifest.json}{source manifest} records all 91 snapshot
hashes. The complete package, including supplementary proofs and exact
arithmetic code, is available at
\href{https://github.com/shjvandekamp/math/tree/main/topics/cylinder-covering}{\texttt{shjvandekamp/math}, topic \texttt{cylinder-covering}}.

\begin{definition}
A cylinder with axis $v\ne0$ and compact perpendicular base $B\subset v^\perp$
is the closed set $B+\R v$. Its cost is the ordinary two-dimensional area
of $B$. The cost of a finite list of cylinders is the sum of its base
areas, with multiplicity. For a convex body $K$, let $\Amin(K)$ be its
minimum orthogonal projection area over all directions.
\end{definition}

Put $H=\sqrt2$ and fix the edge-two regular tetrahedron
\begin{equation}\label{eq:body}
 K=\{(x,y,Hs):0\le s\le1,\ |x|\le1-s,\ |y|\le s\}.
\end{equation}
Its vertices are $(\pm1,0,0)$ and $(0,\pm1,H)$, so every edge has
length two. The two horizontal edge directions are the $x$- and
$y$-axes. Angles below are measured from the corresponding unoriented
horizontal edge line.

\begin{lemma}[Projection normalization]\label{lem:normalization}
The tetrahedron in \eqref{eq:body} satisfies $\Amin(K)=H$.
\end{lemma}
\begin{proof}
The four outward face area vectors are $(0,\pm H,-1)$ and
$(\pm H,0,1)$. For a unit vector $v$, the elementary polytope projection
formula therefore gives
\[
 A(v)=\max(H|v_x|,|v_z|)+\max(H|v_y|,|v_z|).
\]
Indeed, half the sum of the absolute scalar products of the face area
vectors with $v$ counts the projected area: almost every line parallel to
$v$ through the projection meets two faces, and projected edges have area
zero. Set $U=\max(|v_x|,|v_z|/H)$ and
$V=\max(|v_y|,|v_z|/H)$. Then
\[
 (U+V)^2\ge v_x^2+v_y^2+2v_z^2/H^2=1.
\]
Thus $A(v)=H(U+V)\ge H$, with equality for the $x$-axis.
\end{proof}

Throughout the note an area ratio means total cylinder cost divided by
$H$. Define the polynomials
\begin{equation}\label{eq:polynomials}
\begin{split}
 G(q)&=\frac{1+q^2}{2},\qquad D(q)=\frac{1-q^2}{2},
 \qquad h(q)=q^2G(q),\\
 D_0(q)&=G(q)^2+2q^2,\qquad
 C_e(q)=1+e^2h(q)+\frac9{16}e^4,
 \qquad M_*=1+e^2+\frac9{16}e^4.
\end{split}
\end{equation}
In particular, $1\le C_e(q)\le M_*$ on $[-1,1]$.

\begin{theorem}[Effective triangular covers]\label{thm:main}
For every integer $N\ge1$ and $0<e<1$, the explicit construction in
Section~\ref{sec:construction} gives exactly $2N$ labeled cylinders with
compact, nondegenerate triangular perpendicular bases covering every point
of the closed tetrahedron $K$. On the equal grid, their maximum axis-angle
tangent is
\begin{equation}\label{eq:angle-main}
 \sqrt3\,e(1-1/N).
\end{equation}
Write $\Rsec(e,N)$ for their normalized cost and put
\begin{equation}\label{eq:J}
 J(e)=\frac14\int_{-1}^1
       \frac{C_e(q)^2}{\sqrt{1+e^2D_0(q)}}\,dq.
\end{equation}
For $0<e\le1/10$ the following hold:
\begin{align}
 |\Rsec(e,N)-J(e)|&\le\frac{2e}{N},\label{eq:main-error}\\
 N\ge\ceil{480/e}\quad&\Longrightarrow\quad
 \Rsec(e,N)\le\frac12-\frac{e^2}{240},\label{eq:main-count}\\
 N\ge\ceil{2/e^2}\quad&\Longrightarrow\quad
 \Rsec(e,N)\le\frac12-\frac{e^2}{60}
                         +\frac{e^3}{2}+\frac92e^4.
 \label{eq:main-schedule}
\end{align}
Consequently, the count $2N=2\ceil{2/e^2}$ gives
\begin{equation}\label{eq:main-saving}
 \Rsec(e,N)\le\frac12-\frac{73}{15000}e^2<\frac12
 \qquad(0<e\le1/50).
\end{equation}
On this specified count schedule the actual asymptotic cost is
\begin{equation}\label{eq:main-asymptotic}
 \Rsec(e,\ceil{2/e^2})
 =\frac12-\frac{e^2}{60}+\frac{e^3}{2}+O(e^4).
\end{equation}
Finally, at $e=1/10$ and $N=10000$, exact rational bounds give
\begin{equation}\label{eq:benchmark}
\begin{split}
 0.499921585997001674
 &\le\Rsec(1/10,10000)\\
 &\le0.499921585997001774<0.499922.
\end{split}
\end{equation}
\end{theorem}

The coverage assertion is proved for arbitrary positive-width partitions,
with local padding. The quantitative estimates above concern the equal grid;
a nonuniform version of \eqref{eq:main-error} is given below. The coefficient
$1/2$ in \eqref{eq:main-asymptotic} describes this recipe, not an optimum
over all cylinder covers or all compatible sector constructions.

\section{The explicit finite construction}\label{sec:construction}

Let $\mathcal P$ be a partition of $[-1,1]$ into $N$ closed intervals
$I_i=[l_i,r_i]$ of positive width, in their natural order, with disjoint
interiors. Set
\begin{equation}\label{eq:cell-data}
 d_i=r_i-l_i,\quad a_i=d_i/2,\quad c_i=(l_i+r_i)/2,
 \quad A_i=\frac{1+l_ir_i}{2}=G(c_i)-\frac{a_i^2}{2}.
\end{equation}
For $0<e<1$, define
\begin{equation}\label{eq:local-padding}
 K_0=M_*(e+e^2/2)+2e^2,\qquad
 \delta_i=\frac{K_0a_i+(M_*e^2/2)a_i^2}{1-e},\qquad
 M_i=\max_{q\in I_i}C_e(q)+\delta_i.
\end{equation}
The maximum is explicitly attained at the endpoint of largest absolute
value, since $h$ is even and increasing as a function of $|q|$.

Take the intercept triangles
\begin{equation}\label{eq:intercepts}
\begin{split}
 P_i&=\{(0,qt,Ht):q\in I_i,\ 0\le t\le M_i/2\},\\
 Q_i&=\{(pt,0,H(1-t)):p\in I_i,\ 0\le t\le M_i/2\},
\end{split}
\end{equation}
and respectively the axes
\begin{equation}\label{eq:axes}
 v_i=(1,eA_i,Hec_i),\qquad w_i=(-eA_i,1,Hec_i).
\end{equation}
The actual perpendicular bases are the orthogonal projections of $P_i$
onto $v_i^\perp$ and of $Q_i$ onto $w_i^\perp$. The first coordinate of
$v_i$ and the second coordinate of $w_i$ are one. Projection is therefore
injective on the corresponding intercept plane, and neither positive-area
triangle degenerates. Projection also changes every point by a multiple
of the axis, so it leaves the associated cylinder unchanged. This defines
the promised list of $2N$ compact triangular cylinders.

For the equal grid $d=2/N$, every $\delta_i$ equals
\begin{equation}\label{eq:uniform-padding}
 \delta(e,d)=
 \frac{d\{M_*(e+e^2/2)+2e^2+M_*e^2d/4\}}{2(1-e)}.
\end{equation}
In the angular paper's physical notation, $\tau=2e$,
$\alpha_i=A_i/2$, and $\beta_i=c_i/2$. Thus \eqref{eq:axes} are
exactly its secant axes, rather than new directions.

\begin{proposition}[Exact cost and finite angle]\label{prop:cost}
The normalized cost of the construction is
\begin{equation}\label{eq:exact-cost}
 \Rsec(e;\mathcal P)=\frac14\sum_i
       \frac{d_iM_i^2}{\sqrt{1+e^2(A_i^2+2c_i^2)}}.
\end{equation}
For the equal grid its maximum axis-angle tangent is
\eqref{eq:angle-main}, including the case $N=1$.
\end{proposition}
\begin{proof}
Each intercept triangle has area $Hd_iM_i^2/8$. Orthogonal projection
multiplies this by the absolute scalar product of the intercept-plane
unit normal with the unit axis, namely
$(1+e^2(A_i^2+2c_i^2))^{-1/2}$. Summing the two families and dividing
by $H$ proves \eqref{eq:exact-cost}.

The squared tangent from the corresponding horizontal edge is
$e^2(A_i^2+2c_i^2)$. On an equal grid this increases with $|c_i|$:
writing $z=c_i^2$, its factor is
$((1+z)/2-a_i^2/2)^2+2z$, whose derivative is $A_i+2>0$.
At either outer interval, $A_i=|c_i|=1-1/N$. This proves the formula.
For $N=1$, both $A_i$ and $c_i$ vanish, so the directions are the
horizontal equality directions themselves.
\end{proof}

\section{Closed-body selection and the overlap certificate}

We first establish angular selection at every point, without assuming a
monotone root map and without dividing by a possibly zero radius. This is
the shared-boundary first-crossing argument of~\cite{Angular}, included
here with its endpoint cases.

\begin{lemma}[First crossing on a finite partition]\label{lem:crossing}
For every $(x,y,Hs)\in K$ there are indices $i,j$ and nonnegative
numbers $t_0,s_0$ such that
\begin{align*}
 t_0&=s-exc_i,& y_0&=y-exA_i,& l_it_0&\le y_0\le r_it_0,\\
 s_0&=1-s+eyc_j,& x_0&=x+eyA_j,& l_js_0&\le x_0\le r_js_0.
\end{align*}
If either radius is zero, a cylinder of the construction already covers
the point. Otherwise, putting
\[
 \rho=2t_0,\quad \sigma=2s_0,\quad q=y_0/t_0\in I_i,
 \quad p=x_0/s_0\in I_j,\quad X=2x,\quad Y=2y,
\]
gives
\begin{equation}\label{eq:coupled}
 Y=\rho q+eXA_i,\qquad X=\sigma p-eYA_j,
 \qquad \rho+\sigma=2+e(Yc_j-Xc_i).
\end{equation}
\end{lemma}
\begin{proof}
Write $q_0=-1<q_1<\cdots<q_N=1$ for the partition nodes. For
$q\in I_i$, let $L_i(q)=A_i-qc_i$. Direct substitution at either
endpoint gives $L_i(q)=D(q)$. The sector boundaries consequently match
at common nodes. Moreover, $L_i$ interpolates the two nonnegative endpoint
values of $D$, hence $L_i(q)\ge0$ throughout its interval.

At the nodes put $F_k=sq_k+exD(q_k)$. Since
$F_0=-s\le y\le s=F_N$, there is a first $k\ge1$ for which
$F_k\ge y$. If $k=1$, $F_0\le y$ already; if $k>1$, minimality
gives $F_{k-1}<y$. Thus in both cases
$F_{k-1}\le y\le F_k$. Choose $I_i=[q_{k-1},q_k]$. Replacing
$D(q)$ at these endpoints by $A_i-qc_i$ gives exactly
\[
 l_i(s-exc_i)\le y-exA_i\le r_i(s-exc_i).
\]
The outer inequality and $r_i-l_i>0$ imply $t_0=s-exc_i\ge0$.
This also handles $y=F_0$; no monotonicity of the remaining node values
has been used.

For the upper family repeat the argument with
$E_k=(1-s)q_k-eyD(q_k)$, whose endpoint values bracket $x$ because
$|x|\le1-s$. It produces the second row of inequalities and
$s_0=1-s+eyc_j\ge0$.

If $t_0=0$, the bracketing forces $y_0=0$, and the original point is
$xv_i$, on the cylinder through the tip $(0,0,0)$ of $P_i$. If
$s_0=0$, then $x_0=0$, and the point is
$(0,0,H)+yw_j$, on the cylinder through the tip of $Q_j$. In the
remaining case both divisions in the statement are legitimate, and
rearranging the definitions gives \eqref{eq:coupled}.
Every step uses closed inequalities, so this proof covers angular nodes,
all faces and edges, and all four vertices.
\end{proof}

For the positive-radius case of Lemma~\ref{lem:crossing}, put
\begin{equation}\label{eq:weights}
\begin{split}
 \Lambda&=1+e^2A_iA_j,\\
 w_\rho&=1-eqc_j+e^2A_j(A_i-qc_i),\\
 w_\sigma&=1+epc_i+e^2A_i(A_j-pc_j).
\end{split}
\end{equation}
The cell data satisfy $0\le A_i,A_j\le1$, and the interpolation just
proved gives $L_i,L_j\ge0$. Therefore
\begin{equation}\label{eq:weight-positivity}
 w_\rho,w_\sigma\ge1-e>0.
\end{equation}
Solving the first two equations in \eqref{eq:coupled} yields
\[
 X=\frac{\sigma p-e\rho qA_j}{\Lambda},\qquad
 Y=\frac{\rho q+e\sigma pA_i}{\Lambda}.
\]
Substituting these into its third equation gives the exact elimination
identity
\begin{equation}\label{eq:weighted-radii}
 w_\rho\rho+w_\sigma\sigma=2\Lambda.
\end{equation}

\begin{lemma}[Quartic certificate in physical coordinates]\label{lem:square}
For $p,q\in[-1,1]$, put
\[
 N_0=1+e^2G(p)G(q),\quad
 A_-=1-epq+e^2G(p)D(q),\quad
 B_+=1+epq+e^2G(q)D(p).
\]
Then
\begin{equation}\label{eq:square}
\begin{split}
 A_-C_e(q)+B_+C_e(p)-2N_0
 ={}&\frac{e^2}{2}
 \left[p^2-q^2+\frac{epq(1+p^2+q^2)}2\right]^2\\
 &+e^4G(p)G(q)\{p^2D(p)+q^2D(q)\}\\
 &+\frac{e^4}{8}\{9-p^2q^2(1+p^2+q^2)^2\}\\
 &+\frac9{16}e^6\{G(p)D(q)+G(q)D(p)\}\ge0.
\end{split}
\end{equation}
\end{lemma}
\begin{proof}
Write $\Delta=p^2-q^2$ and $Z=1+p^2+q^2$. Expanding the left
side before completing the square gives
\[
 \frac{e^2\Delta^2}{2}
 +e^3pq\{h(p)-h(q)\}
 +e^4G(p)G(q)\{p^2D(p)+q^2D(q)\}
 +\frac9{16}e^4(A_-+B_+).
\]
Here $h(p)-h(q)=\Delta Z/2$ and
$A_-+B_+=2+e^2\{G(p)D(q)+G(q)D(p)\}$. Completing the square
proves the identity. Every displayed term on the right of
\eqref{eq:square} is nonnegative. For the only factor requiring a
bound, $p^2q^2\le1$ and $Z\le3$ already suffice; the exact
nonnegative decomposition is
\[
 9-p^2q^2Z^2
 =9\{(1-p^2)+p^2(1-q^2)\}
  +p^2q^2\{(1-p^2)+(1-q^2)\}(3+Z).
\]
The plus sign inside the square is the sign for the physical coordinates
\eqref{eq:body} and \eqref{eq:coupled}.
\end{proof}

\begin{lemma}[Finite secant perturbation]\label{lem:perturbation}
For the selected intervals and parameters above,
\begin{equation}\label{eq:perturbation}
 w_\rho C_e(q)+w_\sigma C_e(p)-2\Lambda
 \ge-K_0(a_i+a_j)-\frac{M_*e^2}{2}(a_i^2+a_j^2).
\end{equation}
Consequently,
\begin{equation}\label{eq:buffered-certificate}
 w_\rho\{C_e(q)+\delta_i\}
 +w_\sigma\{C_e(p)+\delta_j\}\ge2\Lambda.
\end{equation}
\end{lemma}
\begin{proof}
Writing $b=q-c_i$, define $E_i(q)=(a_i^2-b^2)/2$. Elementary
secant identities give
\begin{equation}\label{eq:secant-errors}
 L_i(q)=D(q)-E_i(q),\qquad 0\le E_i(q)\le a_i^2/2,
 \qquad |A_i-G(q)|\le a_i.
\end{equation}
For the last bound,
\[
 |A_i-G(q)|=\left|-c_ib-\frac{a_i^2+b^2}{2}\right|
 \le|c_i|a_i+a_i^2\le a_i,
\]
using $|c_i|+a_i\le1$. The same statements hold for $p\in I_j$.
Since $D\le1/2$ and $G,A_i,A_j\le1$, they imply
\begin{align}
 |w_\rho-A_-|&\le ea_j+\frac{e^2}{2}(a_j+a_i^2),
 \label{eq:weight-error1}\\
 |w_\sigma-B_+|&\le ea_i+\frac{e^2}{2}(a_i+a_j^2),
 \label{eq:weight-error2}\\
 |\Lambda-N_0|&\le e^2(a_i+a_j).
 \label{eq:denominator-error}
\end{align}
For example, the quadratic term in the first estimate is bounded by
expanding
\[
 A_jL_i(q)-G(p)D(q)
 =[A_j-G(p)]D(q)-A_jE_i(q).
\]
The linear term is bounded by $e|q|\,|c_j-p|\le ea_j$.
The second estimate is identical with the roles exchanged, and the last
follows from
$A_iA_j-G(q)G(p)=(A_i-G(q))A_j+G(q)(A_j-G(p))$.

Multiply \eqref{eq:weight-error1} and \eqref{eq:weight-error2} by
the respective caps, use $C_e\le M_*$, and account for twice the
denominator error. Lemma~\ref{lem:square} then gives
\eqref{eq:perturbation}, with exactly the $K_0$ in
\eqref{eq:local-padding}. Finally,
\[
 w_\rho\delta_i+w_\sigma\delta_j
 \ge(1-e)(\delta_i+\delta_j)
 =K_0(a_i+a_j)+\frac{M_*e^2}{2}(a_i^2+a_j^2).
\]
Adding this to \eqref{eq:perturbation} proves
\eqref{eq:buffered-certificate}.
\end{proof}

\begin{proof}[Proof of the coverage assertion in Theorem~\ref{thm:main}]
Choose the two selected intervals using Lemma~\ref{lem:crossing}. Its
zero-radius cases already give coverage. Otherwise, the point belongs to
the lower selected cylinder whenever $\rho\le M_i$, and to the upper
selected cylinder whenever $\sigma\le M_j$, by
\eqref{eq:intercepts}. If both cylinders missed it, then
$\rho>M_i\ge C_e(q)+\delta_i$ and
$\sigma>M_j\ge C_e(p)+\delta_j$. The strictly positive weights
in \eqref{eq:weight-positivity} would give
\[
 w_\rho\rho+w_\sigma\sigma
 > w_\rho\{C_e(q)+\delta_i\}
  +w_\sigma\{C_e(p)+\delta_j\}\ge2\Lambda,
\]
contradicting \eqref{eq:weighted-radii}. Hence one of the two closed
cylinders covers the point. Equality at a cap boundary is included.
The argument applies to every point of $K$ and every positive-width
partition, without an exceptional boundary set.
\end{proof}

\begin{remark}[Compatible caps beyond constant triangles]\label{rem:caps}
The same proof applies to a cellwise cap $R_i(q)$ whenever
$R_i(q)\ge C_e(q)+\delta_i$. Replace $0\le t\le M_i/2$ in
\eqref{eq:intercepts} by $0\le t\le R_i(q)/2$. If the resulting
intercept sets are compact, their projections are compact cylinder bases;
they need not be triangles. Constant maxima were used above to obtain
triangular bases and the simple exact area formula.
\end{remark}

\section{The limiting integral with an explicit remainder}\label{sec:integral}

The next elementary calculation supplies every constant needed for the
finite area estimates. It is included so that no asymptotic assertion is
used as an effective error bound.

\begin{lemma}\label{lem:integral}
For $0<e\le1/10$,
\begin{equation}\label{eq:J-upper}
 J(e)\le\frac12-\frac{e^2}{60}+\frac{929}{1260}e^4+4e^6.
\end{equation}
Moreover,
$J(e)=1/2-e^2/60+(929/1260)e^4+O(e^6)$ as $e\downarrow0$.
\end{lemma}
\begin{proof}
Let $u=e^2$ and $b=9/16$. Taylor's theorem, using the negative third
derivative, gives for every $z\ge0$
\[
 (1+z)^{-1/2}\le1-z/2+3z^2/8.
\]
Multiply this by
\[
 C_e(q)^2=1+2h(q)u+(h(q)^2+2b)u^2
                 +2bh(q)u^3+b^2u^4.
\]
The constant, linear, and quadratic coefficients of the product integrate
exactly to the first three terms of \eqref{eq:J-upper}. Explicitly,
\[
 \int_{-1}^1h(q)\,dq=\frac8{15},\qquad
 \int_{-1}^1D_0(q)\,dq=\frac{34}{15},
\]
so $\frac14\int(2h-D_0/2)=-1/60$. Also
\[
 h^2-hD_0+\frac38D_0^2
 =\frac{3+44q^2+162q^4-52q^6+19q^8}{128},
\]
whose integral divided by two is $881/2520$. Adding the constant
contribution $2b=9/8$ gives $929/630$ for $2J$, or $929/1260$
for $J$.

The remaining four coefficients of the upper polynomial are
\begin{align*}
 c_3&=\tfrac34hD_0^2-\tfrac12D_0(h^2+2b)+2bh,\\
 c_4&=\tfrac38D_0^2(h^2+2b)-bD_0h+b^2,\\
 c_5&=\tfrac34bhD_0^2-\tfrac12b^2D_0,\qquad
 c_6=\tfrac38b^2D_0^2.
\end{align*}
Using $0\le h\le1$, $0\le D_0\le3$, and dropping the negative
terms gives respectively the pointwise upper bounds
\[
 \frac{63}{8},\quad\frac{1917}{256},\quad
 \frac{243}{64},\quad\frac{2187}{2048}.
\]
For $u\le1/100$,
\[
 \frac{63}{8}+\frac{1917}{25600}
 +\frac{243}{640000}+\frac{2187}{2048000000}<8.
\]
The factor $\frac14\int_{-1}^1$ therefore bounds the full remaining
contribution by $4u^3$, proving \eqref{eq:J-upper}. Taylor's remainder
is also uniformly bounded on $q\in[-1,1]$ and $e\le1/10$, since
$0\le e^2D_0(q)\le3/100$. The same coefficient computation then
gives the stated two-sided asymptotic expansion.
\end{proof}

\section{Finite error, an explicit count, and the cubic term}

In this section $0<e\le1/10$, the grid is equal, and $d=2/N$.
Write $\widehat M_i=\max_{I_i}C_e$, $\delta=\delta(e,d)$, and
\[
 g(q)=(1+e^2D_0(q))^{-1/2},\qquad
 g_i^{\rm sec}=(1+e^2(A_i^2+2c_i^2))^{-1/2}.
\]

\begin{lemma}[All discretization errors]\label{lem:finite-error}
For every $N\ge1$,
\begin{equation}\label{eq:error-upper}
\begin{split}
 \Rsec(e,N)-J(e)\le{}&
 \left(M_*+\frac38M_*^2\right)e^2d
 +M_*\delta+\frac{\delta^2}{2}
 +\frac{M_*^2e^2d^2}{16},
\end{split}
\end{equation}
and
\begin{equation}\label{eq:error-lower}
 \Rsec(e,N)-J(e)\ge-\frac38M_*^2e^2d.
\end{equation}
\end{lemma}
\begin{proof}
There are four changes from the integral to the exact sum.
First, $C_e$ has total variation $2e^2$ on $[-1,1]$. Thus
\[
 \sum_i\int_{I_i}(\widehat M_i-C_e(q))\,dq\le2e^2d.
\]
Since $\widehat M_i+C_e(q)\le2M_*$ and $g\le1$, replacing the
cap by its cell maximum increases the normalized integral by at most
$M_*e^2d$.

Second, $|D_0'(q)|\le6$ and hence $|g'(q)|\le3e^2$. Because
\[
 \sum_i\int_{I_i}|q-c_i|\,dq=d/2,
\]
replacing $g(q)$ by $g(c_i)$ changes the integral, in absolute value,
by at most $3M_*^2e^2d/8$.

Third, with $a_i=d/2$,
\[
 D_0(c_i)-(A_i^2+2c_i^2)
 =G(c_i)a_i^2-a_i^4/4\in[0,a_i^2].
\]
The derivative of $(1+e^2z)^{-1/2}$ in $z$ has absolute value at
most $e^2/2$. Consequently
\[
 0\le g_i^{\rm sec}-g(c_i)\le e^2d^2/8.
\]
Applying this to the unpadded maximum caps gives a nonnegative area
increase at most $M_*^2e^2d^2/16$.

Finally, adding the padding changes their normalized cost by
\[
 \frac14\sum_i d(2\widehat M_i\delta+\delta^2)g_i^{\rm sec}
 \le M_*\delta+\delta^2/2.
\]
The four estimates prove \eqref{eq:error-upper}. Every change except
the midpoint-weight substitution is nonnegative, which also proves
\eqref{eq:error-lower}.
\end{proof}

\begin{proposition}[Uniform error and count bound]\label{prop:uniform}
Equations \eqref{eq:main-error} and \eqref{eq:main-count} hold. More
generally, for an arbitrary positive-width partition with the local
padding \eqref{eq:local-padding},
\begin{equation}\label{eq:nonuniform-error}
 |\Rsec(e;\mathcal P)-J(e)|\le e\max_i d_i.
\end{equation}
\end{proposition}
\begin{proof}
The elementary bounds
\begin{equation}\label{eq:numerical-bounds}
 M_*<\frac{1011}{1000},\qquad M_*^2<\frac{1023}{1000},
 \qquad M_*+\frac38M_*^2<\frac75
\end{equation}
hold throughout $e\le1/10$. We also have
\begin{equation}\label{eq:padding-coarse}
 \delta\le\frac34ed\qquad(0<d\le2).
\end{equation}
For an explicit verification of this last bound, subtract $ed/2$ in
\eqref{eq:uniform-padding} and clear its denominator to obtain
\[
 \delta-\frac{ed}{2}
 =\frac{ed\{(M_*-1)+e(M_*/2+3+M_*d/4)\}}{2(1-e)}.
\]
Here $(M_*-1)/e=e+9e^3/16$ is increasing. Substituting
$e=1/10$, $d=2$ into the remaining upper coefficient gives
\[
 \frac{(M_*(1/10)-1)/(1/10)+M_*(1/10)+3}{2(1-1/10)}
 <\frac52.
\]
Thus $0\le\delta-ed/2\le(5/2)e^2d\le ed/4$.

Apply \eqref{eq:numerical-bounds}, \eqref{eq:padding-coarse},
$e\le1/10$, and $ed\le1/5$ to \eqref{eq:error-upper}. Its entire
right side is at most
\[
 \left(\frac7{50}+\frac{3033}{4000}
       +\frac9{160}+\frac{1023}{80000}\right)ed
 =\frac{77383}{80000}ed<ed.
\]
The absolute lower error in \eqref{eq:error-lower} is smaller than
$ed$ as well. This proves \eqref{eq:main-error}.

For a nonuniform partition set $d=\max_i d_i$. The same maximum-cap
variation estimate is bounded by $2e^2d$. The midpoint estimate uses
$\sum_i d_i^2\le2d$, and the secant estimate uses
$\sum_i d_i^3\le2d^2$. Each local $\delta_i$ is at most
$\delta(e,d)$, by its explicit positive-coefficient formula. The preceding
proof therefore proves \eqref{eq:nonuniform-error} too.

Finally,
\[
 \frac{929}{1260\cdot100}+\frac4{10000}<\frac1{120}.
\]
Lemma~\ref{lem:integral} gives $J(e)\le1/2-e^2/120$.
If $N\ge480/e$, the finite error $2e/N$ is at most $e^2/240$,
which proves \eqref{eq:main-count}.
\end{proof}

\begin{proposition}[The source count schedule]\label{prop:schedule}
Equations \eqref{eq:main-schedule}, \eqref{eq:main-saving}, and
\eqref{eq:main-asymptotic} hold.
\end{proposition}
\begin{proof}
If $N\ge\ceil{2/e^2}$ then $d\le e^2$. This improves the padding
bound to
\begin{equation}\label{eq:padding-fine}
 \delta\le\frac12ed+\frac{21}{10}e^2d
          \le\frac{71}{100}e^3.
\end{equation}
To verify it, use $M_*\le1+(81/80)e^2$, $d\le e^2$, and
$e\le1/10$ to get
\[
 M_*(1+e/2+ed/4)+2e\le1+\frac52e+\frac{11}{10}e^2.
\]
For example the higher-order coefficient is bounded by
\[
 \frac{81}{80}+
 \left(\frac14+\frac{81}{160}\right)\frac1{10}
 +\frac{81}{320\cdot1000}<\frac{11}{10}.
\]
The exact identity
\[
 (1-e)(1+21e/5)-(1+5e/2+11e^2/10)
 =7e/10-53e^2/10\ge0
\]
now proves \eqref{eq:padding-fine} after division by $2(1-e)$.

All remaining terms can be bounded by rational constants. In
\eqref{eq:error-upper} use
\begin{align*}
 \left(M_*+\frac38M_*^2\right)e^2d&\le\frac75e^4,&
 M_*\delta&\le\frac12e^3+\frac{109}{50}e^4,\\
 \delta^2/2&\le\frac3{1000}e^4,&
 M_*^2e^2d^2/16&\le\frac1{1000}e^4.
\end{align*}
For the second inequality, \eqref{eq:padding-fine} supplies the
$e^3/2+(21/10)e^4$ part, and
\[
 (M_*-1)\delta
 \le\frac{81}{80}e^2\frac{71}{100}e^3
 <\frac2{25}e^4.
\]
The others follow directly from \eqref{eq:numerical-bounds},
\eqref{eq:padding-fine}, and $e^2\le1/100$. With
$4e^6\le e^4/25$, the total quartic coefficient from
Lemma~\ref{lem:integral} and these inequalities is
\[
 \frac{929}{1260}+\frac75+\frac{109}{50}
 +\frac3{1000}+\frac1{1000}+\frac1{25}<\frac92.
\]
This proves \eqref{eq:main-schedule}. If $e\le1/50$, then
\[
 \frac1{60}-\frac e2-\frac92e^2
 \ge\frac1{60}-\frac1{100}-\frac9{5000}
 =\frac{73}{15000},
\]
giving \eqref{eq:main-saving}.

For the exact schedule $N=\ceil{2/e^2}$, one has
$d=e^2+O(e^4)$. Formula \eqref{eq:uniform-padding} gives
\[
 \delta=ed/2+(7/4)e^2d+O(e^3d+e^2d^2)
       =e^3/2+O(e^4).
\]
By the first three estimates in Lemma~\ref{lem:finite-error}, the
unpadded maximum-cap secant sum differs from $J(e)$ by $O(e^2d)$.
The coefficient of its linear padding term is
\[
 \frac12\sum_i d\widehat M_i g_i^{\rm sec}=1+O(e^2),
\]
since $\widehat M_i=1+O(e^2)$, $g_i^{\rm sec}=1+O(e^2)$
uniformly, and $\sum_i d=2$. The quadratic padding term is
$O(\delta^2)$. Together with Lemma~\ref{lem:integral}, these facts
prove \eqref{eq:main-asymptotic}.
\end{proof}

In physical angular notation $\tau=2e$, the schedule bound becomes
\[
 \Rsec\le\frac12-\frac{\tau^2}{240}
                   +\frac{\tau^3}{16}+\frac9{32}\tau^4,
 \qquad N\ge\ceil{8/\tau^2}.
\]
The coefficient conversions matter when comparing statements written in
different source coordinate systems.

\section{An exact finite example and its arithmetic certificate}
\label{sec:certificate}

Set $e=1/10$, $N=10000$, and $d=1/5000$. Then
\[
 M_*=1.01005625,\qquad
 \delta=0.00001400626791979166\ldots<0.00001401.
\]
The definition of $\delta$ is the exact rational expression
\eqref{eq:uniform-padding}, not a decimal truncation. Its maximum axis
tangent is $(\sqrt3/10)(9999/10000)$.

Even without computing the finite sum, the estimates already proved give
\begin{equation}\label{eq:analytic-benchmark}
 \Rsec(1/10,10000)<0.499929.
\end{equation}
Indeed, the upper bound for $J(e)$ in \eqref{eq:J-upper} is less than
$0.499911064$, the first term of \eqref{eq:error-upper} is less than
$0.0000028$, its last term is less than $0.000000000026$, and
$M_*\delta+\delta^2/2<0.000014164209$. Their sum is below the
claimed value. This analytic estimate is distinct from the tighter
exact-sum enclosure \eqref{eq:benchmark}.

\subsection{Finite rational bounds}
For $m\ge0$, define the rational Taylor polynomial
\[
 T_m(z)=\sum_{k=0}^m(-1)^k\frac{\binom{2k}{k}}{4^k}z^k.
\]
Taylor's integral remainder has the sign of $(-1)^{m+1}$ for
$(1+z)^{-1/2}$ on $z\ge0$. Hence, for every nonnegative $z$,
\begin{equation}\label{eq:taylor-enclosure}
 T_9(z)\le(1+z)^{-1/2}\le T_8(z).
\end{equation}
For rational $e$, put $D_i=A_i^2+2c_i^2$ in the exact finite
construction. The two rational numbers
\begin{equation}\label{eq:rational-certificate}
 L_N=\frac d4\sum_i M_i^2T_9(e^2D_i),\qquad
 U_N=\frac d4\sum_i M_i^2T_8(e^2D_i)
\end{equation}
satisfy $L_N\le\Rsec(e,N)\le U_N$ by
\eqref{eq:exact-cost} and \eqref{eq:taylor-enclosure}.
Thus the numerical enclosure reduces to a finite rational calculation
with specified polynomials and no floating-point acceptance criterion.

There is no need to materialize all $N$ terms. On positive-midpoint cells
the cap maximum is $C_e(c_i+d/2)$, so both the cap and the squared axis
norm are polynomials in the index. For even $N$, the positive midpoints
are $(2j-1)/N$, $1\le j\le N/2$; for odd $N$, they are $2j/N$,
$1\le j\le(N-1)/2$, with a separately included central cell.
The lower polynomial in \eqref{eq:rational-certificate} has degree at
most 44 in $j$, and the upper polynomial has degree at most 40.
Reflection accounts for the negative cells. The exact power sums
$S_r(m)=\sum_{j=1}^m j^r$ can be evaluated recursively from
\[
 S_0(m)=m,\qquad
 S_r(m)=\frac{(m+1)^{r+1}-1-
                \sum_{k=0}^{r-1}\binom{r+1}{k}S_k(m)}{r+1}
 \quad(r\ge1).
\]
This recurrence is simply the telescoping sum of
$(j+1)^{r+1}-j^{r+1}$. It makes \eqref{eq:rational-certificate} a
constant-degree rational certificate whose integer lengths grow with
$\log N$ for fixed rational $e$.

The separate implementation
\repo{verification/secant_verify.py}{\texttt{verification/secant\_verify.py}}
evaluates these bounds with exact fractions. It gives the following
outward-rounded decimal enclosures:
\begin{center}
\small
\begin{tabular}{@{}crl@{}}
\toprule
$e$ & Cylinders $2N$ & Certified normalized area interval\\
\midrule
$1/10$ & 20,000 & $[0.499921585997001674,\ 0.499921585997001774]$\\
$1/100$ & 40,000 & $[0.499998863418315727,\ 0.499998863418315728]$\\
$1/1000$ & 4,000,000 & $[0.499999983836322736,\ 0.499999983836322737]$\\
\bottomrule
\end{tabular}
\end{center}
In particular, the first row proves \eqref{eq:benchmark} and completes
Theorem~\ref{thm:main}. The exact rationals in
\eqref{eq:rational-certificate}, rather than the displayed decimal
precision, are the underlying certificate.

\subsection{What the independent verifier checks}
The acceptance proof reconstructs the secant interpolation, the coupled
elimination, the polynomial identity \eqref{eq:square}, and the exact
perturbation before estimating it. It checks the rational inequalities in
the finite error bounds and evaluates \eqref{eq:rational-certificate}.
It uses the shared exact polynomial and power-sum primitives in
\repo{verification/verify.py}{\texttt{verification/verify.py}}, but does
not import the cylinder generator into those acceptance calculations.
For $N=1,2,3,7$, the resulting sums are also checked against a direct sum
of rational integer-square-root enclosures.

The generator is imported separately for 36 representative output checks
covering both families, $e=1/100,1/10$, and $N=1,2,3,10000$.
These check the axes, intercept vertices, padding, parallel displacement
under projection, perpendicularity, positive Gram determinants, and exact
base areas. Coordinates are stored as $(x,y,z/\sqrt2)$, with metric
$\operatorname{diag}(1,1,2)$; thus this check retains the physical metric
rather than treating the scaled coordinates as Euclidean. These finite
output checks supplement the universal argument in Section 3; they do not
replace it with sampled coverage.

The \repo{secant_construction.py}{indexed constructor} returns any one of
the $2N$ cylinders from its index, using exact rational scaled coordinates
and the exact area expression. For fixed rational $e$ it needs a constant
number of rational values with bit lengths $O(\log N)$, rather than an
explicit list of every cylinder.

\subsection{A larger proved range for a source-form cap}
At $e=1/10$, the source-form compatible cap with
$\tau=1/5$ and $\eta=1/1000$ has twice its physical radial cutoff equal
to
\[
 R_L(q)=1+e^2\{h(q)+1/125\}.
\]
Its excess over the quartic cap is independent of $q$:
\[
 R_L(q)-C_e(q)=\frac{19}{800000}=0.00002375>\delta(e,1/5000).
\]
Remark~\ref{rem:caps} therefore proves a closed-body cover with this
continuous cap and the same 10,000 secant intervals. Its cell-maximum
triangular enlargement also covers. This is a new admissibility proof at
these parameters, rather than an assertion that they met the older
paper's sufficient thresholds. The present triangular cap is smaller
on every cell than that enlarged source-form cap and, with identical
axes, has strictly smaller cost.

\section{Comparison with the identified published ranges}
\label{sec:comparison}

The finite benchmark allows at most 20,000 cylinders. The maximum tangent
is bounded by $(\sqrt3/10)(9999/10000)$ from the corresponding horizontal
equality edges. The radial and generic estimates below allow even the
slightly larger tangent budget $\sqrt3/10$, so also apply at this
benchmark. The comparison concerns the identified source recipes on the
explicit sufficient ranges and count conventions in their proofs. A
sufficient threshold is not treated as a necessary condition for all
geometrically valid covers.

\begin{center}
\small
\renewcommand{\arraystretch}{1.18}
\begin{tabular}{@{}>{\raggedright\arraybackslash}p{0.46\textwidth}>{\raggedright\arraybackslash}p{0.48\textwidth}@{}}
\toprule
Published recipe and convention & Obstruction at the benchmark\\
\midrule
Radial linear, cubic, fractional, and square-root caps on their explicit
source ranges; arbitrary partitions and tags, actual projected bases or
larger enclosures
& $R\ge1/2-1/23328000>0.49999995$, allowing retuned source tilt.\\
Quantitative angular maximum caps, including the retained-mesh sufficient
budget and the cubic-buffer family
& $R\ge1/2-1/93312000$.\\
General compatible angular family on the explicit proof witness
& $R\ge1/2-1/5529600000$; the full list also needs at least 36,864,000 cylinders.\\
Separate direct endpoint and squared constructions
& At least 392,000,002 and 800,000,000 cylinders, respectively.\\
Alternative local Cartesian certificate
& Any resulting complete realization has $R>0.499999999996$.\\
Rational or singular full retained-grid recipe
& $R\ge637/1250=0.5096$.\\
Rational product/logarithmic recipe with its published saving-error
allocation
& More than 20 million selected main cylinders in one family, including
after finite-subcover selection.\\
Nilpotent recipe with its prescribed apex lists
& At least 34,134 cylinders from the apex lists alone.\\
Shortened beams with the prescribed strip
& A strip axis deviates by more than $\arctan(1/3)$, outside the angle
budget even after tetrahedron symmetries.\\
\bottomrule
\end{tabular}
\end{center}
Every listed area floor exceeds the new certified upper bound
$0.499922$. The table summarizes proved comparisons in the public
companion notes: the
\repo{research/published-proof-ranges.md}{explicit proof-range audit},
the \repo{research/published-radial-angular-comparison.md}{radial and
angular actual-area comparison}, and the
\repo{verification/published-generic-comparison.md}{generic realization
and count comparison}. Those notes include the complete derivations and
the precise source locations; their geometric count arguments are not
replaced here by a claim about asymptotic coefficients.

For clarity, two features of the comparison merit spelling out.
First, the general compatible angular theorem uses the explicit
witness~\cite{LeanSources}
\begin{equation}\label{eq:old-general-range}
 E\ge0,\qquad0<\eta\le1/480,\qquad0<\tau\le1/2,
 \qquad(10+4E)\tau\le\eta.
\end{equation}
Its source area estimate is
\[
 \left|R-\left(\frac12+(2\eta-1/240)\tau^2\right)\right|
 \le(6E+2)\tau^4.
\]
The source's extra expansion condition $E\tau^2\le1/16$ follows
from \eqref{eq:old-general-range}. Since
$6E+2\le2(10+4E)$ and $\tau\le1/2$, the positive
$2\eta\tau^2$ absorbs this remainder. Also $\tau\le1/4800$.
This proves the table's corresponding area floor. On each cell the
constant angular coefficient must approximate $q/2$ with error at most
$E\tau^2$ at both endpoints. The width is therefore at most
$4E\tau^2$. Positive-width cells force $E>0$, and their full two-family
count is at least
\[
 \frac1{E\tau^2}\ge\frac{(10+4E)^2}{E\eta^2}
 \ge\frac{160}{\eta^2}\ge36{,}864{,}000,
\]
where $(10+4E)^2-160E=(4E-10)^2\ge0$. The separate direct
endpoint and squared certificates use their own ranges, rather than
being silently assigned \eqref{eq:old-general-range}.

Second, immediate improvements of the older angular proof belong in
the baseline. Retaining its actual mesh $\Delta$ gives the sufficient
condition
\[
 2(1-\tau)\eta\ge\frac74\tau+
                       \frac{\Delta}{4\tau}+\frac\Delta4.
\]
On the same count schedule, $\Delta\sim\tau^2/4$ gives a canonical
$e$-cubic coefficient $29/2$, and $\Delta=o(\tau^2)$ gives coefficient
14. The coarser substitution $\Delta\le\tau^2$ gives coefficient 16.
These are stronger prior-art corollaries than choosing only the larger
printed cubic coefficient. The table's actual-area floor includes this
optimized retained-mesh condition, not just one favorable printed example.

None of these comparisons asserts that every possible witness of an
existential smallness theorem obeys the exhibited numerical witness.
Nor do they cover redesigned tile deletion, different product shapes,
new error allocations, or new admissibility proofs for generic methods.
The cap extension in Section~\ref{sec:certificate} itself demonstrates
why that distinction is necessary. There is no global lower bound for
all tetrahedron cylinder covers in this note.

\section{Supplementary results, reproduction, and verification status}

The repository also contains a different finite realization using frozen
radial sweeps. Its complete proof is in
\repo{THEOREM.md}{\texttt{THEOREM.md}}, with details in the
\repo{research/finite-triangle-optimality.md}{finite-triangle optimality
note}. Within its stated class---symmetric intercept triangles containing
each full constant-maximum-cap frozen sweep---the enclosure is optimal
and its count schedule has sharp cubic coefficient 2. These are
supplementary results whose proofs are in the linked files, not additional
theorems proved in this manuscript. The secant triangles used here do not
contain each cell's full continuously tilted sweep, so that restricted
lower bound does not apply. Sharpness of the constant $9/16$ within the
constant-quartic continuous cap family is likewise proved separately in
\repo{research/cap-analysis.md}{the cap analysis}; global cap optimality
is not asserted.

From the repository root, one command verifies all source pins and both
exact verification outputs:
\begin{quote}\small\ttfamily
python3 topics/cylinder-covering/reproduce.py
\end{quote}
An optional query returns one indexed cylinder:
\begin{quote}\small\ttfamily
python3 topics/cylinder-covering/secant\_construction.py\par
\hspace*{1em}--epsilon 1/10 --n 10000 --index 0
\end{quote}
The two displayed lines form one command. Only the Python standard
library is required. The
\repo{verification/secant-report.md}{independent certificate report} and
\repo{verification/secant-last-run.txt}{saved secant verification output}
document the checked identities, inequalities, and outward intervals.

This package was developed with AI assistance using OpenAI Codex.
The evidence for the new theorem consists of the universal written proof
and the identified exact-arithmetic checks. It is not a claim of human
peer review. The new theorem has not been formalized or kernel checked
in Lean. Pinned Lean files are evidence for the cited prior-art definitions,
proof witnesses, and area formulas; their presence is not a formal
verification of the new quartic/secant theorem.

\begin{thebibliography}{99}
\bibitem{Angular}
OpenAI,
\emph{Finite angular cylinder covers below the half-area bound},
27 September 2026, OpenAI math family 100.
\href{https://github.com/openai/math/tree/1581d4cbdd69a0faf34b2089614fe2a814316a6a/preprints/Finite-angular-cylinder-covers-below-the-half-area-bound-September-27-2026}{Pinned manuscript and source files}.

\bibitem{Radial}
OpenAI,
\emph{Finite triangular approximation of radial sweeps},
27 September 2026, OpenAI math family 100.
\href{https://github.com/openai/math/tree/1581d4cbdd69a0faf34b2089614fe2a814316a6a/preprints/Finite-triangular-approximation-of-radial-sweeps-September-27-2026}{Pinned manuscript and source files}.

\bibitem{Ruled}
OpenAI,
\emph{Finite cylinder approximation of ruled sets},
27 September 2026, OpenAI math family 100.
\href{https://github.com/openai/math/tree/1581d4cbdd69a0faf34b2089614fe2a814316a6a/preprints/Finite-cylinder-approximation-of-ruled-sets-September-27-2026}{Pinned manuscript and source files}.

\bibitem{Slope}
OpenAI,
\emph{Slope-field perturbations of the two-cylinder covering},
27 September 2026, OpenAI math family 100.
\href{https://github.com/openai/math/tree/1581d4cbdd69a0faf34b2089614fe2a814316a6a/preprints/Slope-field-perturbations-of-the-two-cylinder-covering-September-27-2026}{Pinned manuscript and source files}.

\bibitem{LeanSources}
OpenAI,
\emph{Cylinder-covering formal source files},
repository revision\newline
\texttt{1581d4cbdd69a0faf34b2089614fe2a814316a6a}.
The pinned local
\repo{sources/lean/OAI/Geometry/CylinderCovering/Angular/GeneralCoverage.lean}{general coverage},
\repo{sources/lean/OAI/Geometry/CylinderCovering/Angular/GeneralAssembly.lean}{assembly and expansion},
\repo{sources/lean/OAI/Geometry/CylinderCovering/Angular/GeneralDefinitions.lean}{admissibility definitions},
and \repo{sources/lean/OAI/Geometry/CylinderCovering/Angular/GeneralVolume.lean}{physical area identity}
give the explicit compatible-cap witness discussed in
Section~\ref{sec:comparison}; other source locations are identified in the
\repo{research/published-proof-ranges.md}{proof-range audit}.
\end{thebibliography}
\end{document}
